% ============================================================
% 2  Luttinger liquid as a Taylor expansion -> Lifshitz theory
% ============================================================
\begin{frame}[t]
\ft{Luttinger liquid as a Taylor expansion}\par
\vspace{0.4cm}
\centering
{\small$\displaystyle
\mathcal L = i\,\partial_x\theta\,\partial_\tau\phi + \tfrac12(\partial_x\theta)^2
+ \underbrace{\colorbox{cLL!20}{$\mystrut\tfrac{\lambda_2}{2}(\partial_x\phi)^2$}}_{\text{\scriptsize\color{cLL} Luttinger liquid}}
\onslide<2->{+ \underbrace{\colorbox{cQCP!15}{$\mystrut\tfrac12(\partial_x^2\phi)^2 + \tfrac{\lambda_4}{4}(\partial_x\phi)^4$}}_{\text{\scriptsize\color{cQCP} next order}}}
+ \dots$\par}
\vspace{0.1cm}
{\scriptsize density $\rho=\partial_x\phi=S^z$,\qquad $\lambda_2=(2\pi/K_s)^2$\par}
\vspace{0.3cm}
\begin{columns}[T]
\begin{column}{0.52\textwidth}
\begin{itemize}
 \item Expand in gradients of $\phi$: \colorbox{cLL!20}{\mystrut Luttinger liquid} keeps the first term
 \onslide<2->{\item $K_s\to\infty$: $\lambda_2\to0$, keep the \colorbox{cQCP!15}{\mystrut next terms}: Lifshitz theory, $\omega\propto k^2$}
 \onslide<3->{\item $\lambda_4$ relevant: \colorbox{citron!40}{\mystrut interacting}, non-conformal critical point}
\end{itemize}
\end{column}
\begin{column}{0.46\textwidth}
\centering
\onslide<3->{\resizebox{\linewidth}{!}{\phaseline{$\lambda_2$}{}}\\[2pt]
{\scriptsize $z = 2-\tfrac{2}{243}\epsilon^2\approx1.99$ ($\epsilon=2-d=1$)\par}
\vspace{0.15cm}
{\tiny\color{citegray} A.~Nahum, arXiv:2506.21540 (2025)\\
R.~Flores-Calder\'on, M.~Z\"undel, arXiv:2511.09097 (2025)\\
N.~Chepiga, F.~H.~L.~Essler, arXiv:2609.31319 (2026)\par}}
\end{column}
\end{columns}
\onslide<2->{\tk{How does $K_s\to\infty$ at a finite lattice coupling? At a critical point}}
\end{frame}
% Notes: The Luttinger liquid is the first term of a Taylor expansion of the energy in gradients of phi; its coefficient is set by the spin Luttinger parameter, lambda_2 = (2 pi / K_s)^2.
% Notes: Order-parameter correlations in the LL decay as |x|^{-4/K_s}: the larger K_s, the slower. Naively K_s -> infinity means no decay at all.
% Notes: Step 2: but then the first Taylor coefficient vanishes and we must keep the next terms, (d_x^2 phi)^2 and (d_x phi)^4. Note (d_x^2 phi)^2 = (d_x rho)^2: exactly the gradient cost of the density from the previous slide. This is Nahum's Lifshitz field theory, omega ~ k^2.
% Notes: Step 3: tuning lambda_2 through zero: LL -> ferromagnet (S^z condenses). The quartic term is relevant, so the critical point is interacting and not conformal; z is close to 2 (epsilon expansion, numerics). Our question: where does K_s diverge in a real lattice model? -> at this critical point.
